\providecommand{\rank}{\operatorname{r}}
\providecommand{\RR}{\mathbb R}
\providecommand{\CC}{\mathbb C}
\providecommand{\KK}{\mathbb K}

% ---------- Exercise 54 ----------
\setcounter{enumi}{53}
\item {\bfseries\boldmath
State, with justification, whether each of the following assertions is true or false for every square matrix $A$ of order $3$.
}

\begin{enumerate}[label=\alph*)]

\item $\det(I_3+A)=1+\det A$.

\begin{answer}
False. For $A=I_3$, $\det(I_3+A)=\det(2I_3)=8$, whereas $1+\det A=2$.
\end{answer}

\item $\det(A^4)=(\det A)^4$.

\begin{answer}
True. Repeated multiplicativity gives $\det(A^4)=(\det A)^4$.
\end{answer}

\item If $\det A=0$, then the linear system $AX=0$ has infinitely many solutions.

\begin{answer}
True. A homogeneous system is consistent. If $\det A=0$, its echelon
form has fewer than three pivots, hence at least one free unknown. Therefore
$AX=0$ has infinitely many solutions.
\end{answer}

\item If $\det A=0$, then the linear system $AX=\begin{bmatrix}5\\-2\\1\end{bmatrix}$ has infinitely many solutions.

\begin{answer}
False. Take $A=0_{3\times3}$. Its determinant is zero, but
$AX=(5,-2,1)^T$ is inconsistent.
\end{answer}

\item If the reduced echelon form equivalent to $A$ has a zero row, then $\det A=0$.

\begin{answer}
True. A zero row in the reduced form gives $\rank(A)<3$, so $A$ is
singular and $\det A=0$.
\end{answer}
\end{enumerate}
